\section{Formal mechanisms and their exact hypotheses}
\label{sec:formal-mechanisms}

The formal work is most useful when its hypotheses remain visible. A theorem
that identifies an obstruction is different from a theorem that makes that
obstruction vanish. A theorem that reconstructs a matrix from compatible
residues is different from a theorem that attaches a Galois representation to
an automorphic object. This section develops the algebra behind the checked
interfaces, including examples where an attractive stronger statement fails.
The purpose is to make the output of each mechanism clear enough to use in a
future arithmetic argument.

The coefficient rings, groups, and local restrictions below are part of the
input. Their connection to a particular number field or automorphic packet
requires additional theorems. Nothing here constructs a global automorphic
carrier. The matrix results do give actual integral matrices, rather than a
formal symbol for a proposed limit. Their scope is therefore substantial, but
local and explicit.

\subsection{A cyclic cocycle is constrained by the norm}

Let $C_e=\langle\gamma\rangle$ be a finite cyclic group of order $e$, acting
linearly on a module $M$. Write $T$ for the action of $\gamma$ and set
\[
 N=1+T+\cdots+T^{e-1}.
\]
The relation $T^e=1$ implies $N(T-1)=(T-1)N=0$. These identities explain why
the image of $T-1$ lies in the kernel of $N$, but they do not identify the two
spaces.

A one-cocycle is a function $a:C_e\to M$ satisfying
\[
 a(gh)=a(g)+g\cdot a(h).
\]
Putting $g=h=1$ gives $a(1)=0$. If $x=a(\gamma)$, repeated application of the
cocycle equation gives
\[
 a(\gamma^r)=\sum_{j=0}^{r-1}T^j x\qquad(1\le r\le e).
\]
The group relation $\gamma^e=1$ therefore forces $Nx=0$. Conversely, if
$Nx=0$, this formula defines a cocycle. Whenever exponents wrap around $e$,
the discrepancy is a translate of $Nx$, which is zero. Thus evaluating at
the chosen generator identifies cocycles with $\ker N$.

For a zero-cochain $y\in M$, its coboundary at $\gamma$ is $(T-1)y$. Taking
the quotient by these changes gives the precise presentation
\begin{equation}
 H^1(C_e,M)\simeq \ker N/\operatorname{im}(T-1).
 \label{eq:appendix-cyclic-H1}
\end{equation}
If the norm operator is identically zero, this reduces to
$M/\operatorname{im}(T-1)$. That reduction has a hypothesis. Dropping the
kernel from the general formula introduces classes that cannot satisfy the
cyclic relation.

Suppose now that $e$ is invertible in the coefficient ring. Define
$S=\sum_{j=0}^{e-1}jT^j$. A direct coefficient calculation gives
\[
 (T-1)S=e-N.
\]
For $x\in\ker N$, take $y=e^{-1}Sx$. Then $(T-1)y=x$, so every cocycle is
a coboundary and $H^1(C_e,M)=0$. The averaging argument uses division by
$e$ at its decisive step. In residue characteristic dividing $e$, the same
display remains an identity but no longer constructs $y$.

\begin{example}[The same generator action gives different cyclic cohomology]
Over $\mathbb F_3$, let $M=\operatorname{Mat}_2(\mathbb F_3)$ and let $T$
be conjugation by $F=\operatorname{diag}(1,-1)$. The diagonal entries of a
matrix form the $+1$ eigenspace of $T$. The off-diagonal entries form its
$-1$ eigenspace. Both spaces have dimension two.

For $C_2$, the norm $1+T$ is multiplication by $2$ on the diagonal space and
zero on the off-diagonal space. Hence its kernel is the off-diagonal space.
The operator $T-1$ is invertible on that space and zero on the diagonal
space. Formula~\eqref{eq:appendix-cyclic-H1} gives $H^1(C_2,M)=0$.

Let instead the generator of $C_6$ act by the same $T$. On the diagonal
space, its norm is multiplication by $6=0$ in $\mathbb F_3$. On the
off-diagonal space, the alternating six-term sum is also zero. Now $N=0$
on all of $M$, while the image of $T-1$ is still only the off-diagonal
space. Therefore $H^1(C_6,M)$ has dimension two. The generator matrix
alone does not determine this cohomology without the order and coefficient
characteristic.
\end{example}

In Frobenius coordinates the relevant action is often $T(X)=FXF^{-1}$.
The images of $T-1$ and $1-T$ coincide, since multiplication of a source
vector by $-1$ exchanges the signs. This is the bridge to the concrete
quotient by Frobenius coboundaries. The recovered v0.18 source proves that
bridge under the explicit zero-norm hypothesis. It does not claim that every
Frobenius quotient is automatically an arithmetic local cohomology group.
One must first identify the acting group and the coefficient module.

\subsection{A square-zero defect gives a genuine obstruction class}

Let $q:R'\to R$ be a ring map with kernel $I$, and assume $I^2=0$. Let
$\rho:G\to R^\times$ be a group representation. Choose unit-valued lifts
$r_g\in (R')^\times$ with $q(r_g)=\rho(g)$ and $r_1=1$. Such a choice is
an input to this calculation. It is not already a multiplicative lift.
Define its normalized defect by
\begin{equation}
 c(g,h)=(r_gr_h-r_{gh})r_{gh}^{-1}\in I.
 \label{eq:appendix-normalized-defect}
\end{equation}
Equivalently, $r_gr_h=(1+c(g,h))r_{gh}$. The coefficients are elements of
the actual ring kernel, with its additive group, rather than an unrelated
vector space of the same dimension.

The action of $G$ on $I$ is $g\cdot x=r_gxr_g^{-1}$. It is independent of
the chosen unit lift. Indeed, any other unit lift has the form $(1+b)r_g$
with $b\in I$. Since $(1+b)^{-1}=1-b$ and products of two elements of $I$
vanish, conjugating $x\in I$ by $1+b$ leaves $x$ unchanged. The same
observation shows that the defect in the multiplication of the $r_g$ does
not prevent these conjugations from forming a genuine group action on $I$.

Compute $(r_gr_h)r_\ell$ and $r_g(r_hr_\ell)$. After expressing both in
terms of $r_{gh\ell}$, the first is
\[
 (1+c(g,h)+c(gh,\ell))r_{gh\ell},
\]
while the second is
\[
 (1+g\cdot c(h,\ell)+c(g,h\ell))r_{gh\ell}.
\]
All omitted cross products belong to $I^2$. Associativity therefore yields
\begin{equation}
 g\cdot c(h,\ell)-c(gh,\ell)+c(g,h\ell)-c(g,h)=0.
 \label{eq:appendix-cocycle}
\end{equation}
This is the usual inhomogeneous two-cocycle identity. Normalization at the
identity follows from $r_1=1$.

Change the chosen lifts to $r'_g=(1+b(g))r_g$, where $b:G\to I$ and
$b(1)=0$. Expanding without commuting the factors gives
\[
 r'_gr'_h
 =(1+b(g)+g\cdot b(h)+c(g,h))r_{gh}.
\]
To compare with $r'_{gh}$, multiply by $1-b(gh)$. The new defect is
\begin{equation}
 c'(g,h)=c(g,h)+g\cdot b(h)-b(gh)+b(g).
 \label{eq:appendix-defect-change}
\end{equation}
Thus changing the section adds a two-coboundary. The class $[c]\in H^2(G,I)$
does not depend on the section.

\begin{proposition}[Exact lifting criterion]
With the hypotheses above, a unit-valued representation lifting $\rho$
exists if and only if $[c]=0$.
\end{proposition}

\begin{proof}
If $c=\delta b$, replace $r_g$ by $(1-b(g))r_g$. Equation
\eqref{eq:appendix-defect-change} shows that the corrected defect is zero.
The corrected units are therefore multiplicative and reduce to $\rho$.
The equality $c(1,1)=0$ forces $b(1)=0$ for a primitive in the displayed
normalization. Conversely, a multiplicative unit lift has zero defect.
Its difference from the original section is a kernel-valued one-cochain,
so the same change formula proves that the original class is zero.
\end{proof}

The cocycle identity proves compatibility of the defect. It does not prove
that its cohomology class vanishes. This distinction is indispensable when
an obstruction is used as the target of a resolution argument.

\begin{example}[An obstruction that really does not vanish]
Reduce $\mathbb Z/27\mathbb Z$ to $\mathbb Z/9\mathbb Z$. Its kernel is
$9\mathbb Z/27\mathbb Z$, whose square is zero. Map a generator of $C_3$
to the unit $4$ modulo $9$. This defines a representation because
$4^3=1$ modulo $9$.

A lift of the generator must be $4+9t$ modulo $27$, for $t=0,1,2$.
Direct expansion gives
\[
 (4+9t)^3\equiv 4^3\equiv 10\pmod {27}.
\]
None has cube $1$. Hence no representation lift exists and the obstruction
class is nonzero. Conjugation is trivial in this scalar example. The
kernel is additively $\mathbb F_3$, and the standard cyclic resolution gives
$H^2(C_3,\mathbb F_3)=\mathbb F_3$ for the trivial action. This is a
failure of lifting, rather than a failure to find a convenient section.

The defect can also be seen without an existence argument. Choose the
section $r_1=1$, $r_\gamma=4$, and $r_{\gamma^2}=16$ modulo $27$. These
reduce to the required units modulo $9$. The products of consecutive
powers have no defect until the exponent passes through three. At that
point $4^3=10=1+9$ modulo $27$. Consequently the normalized defect is
$9$ when the exponents wrap around and zero otherwise. Under the
identification $I\simeq\mathbb F_3$, this is the ordinary carry cocycle.

For a proposed correction $b$, sum the equations
$c(\gamma^i,\gamma)=b(\gamma)-b(\gamma^{i+1})+b(\gamma^i)$ for
$i=0,1,2$. The right side telescopes to $3b(\gamma)=0$, using
$b(1)=0$ and the trivial action. The left side is $9\ne0$ in $I$.
There is therefore no primitive. This explicit calculation recovers the
nonzero obstruction class directly from its normalized representative.
\end{example}

For topological groups the cochains and correcting primitives must belong
to the required continuous complex. An abstract primitive need not be
continuous. The recovered continuous modules impose continuous sections and
local restrictions, with a discrete coefficient ring. They preserve those
hypotheses through their lifting criteria. Replacing continuous cohomology
by abstract group cohomology would require a separate comparison.

\subsection{Local conjugacy leaves a second gluing obstruction}

Suppose an absolute two-cocycle has been cancelled by a continuous
one-cochain $b_0$. Prescribed local lifts need not agree with the resulting
global lift. Their remaining differences are local one-cocycles. Changing
$b_0$ to another absolute primitive changes these differences by the
restriction of a global one-cocycle. Conjugating the local lift at index
$i$ by a unit $1+x_i$, with $x_i\in I$, changes the local cocycle by a
local coboundary. The conjugators $x_i$ are independently chosen at each
local index. Requiring one common conjugator is a stronger problem.

After taking the local coboundary quotients, the remaining class belongs to
\begin{equation}
 \operatorname{coker}\left(
 H^1_{\mathrm{cont}}(G,I)\longrightarrow
 \prod_i H^1_{\mathrm{cont}}(G_i,I)\right).
 \label{eq:appendix-gauge-cokernel}
\end{equation}
Here the local groups, their maps into $G$, and the precise continuous
coefficient action must be supplied. The recovered source constructs the
corresponding quotients directly in its square-zero data.

To see the criterion, let $a_i$ be the local mismatch after $b_0$. A match
up to the allowed local conjugations exists exactly when
\[
 a_i=\operatorname{res}_i(z)+\delta x_i
\]
for a global continuous one-cocycle $z$ and kernel elements $x_i$. This is
exactly the assertion that the image of $(a_i)_i$ in
\eqref{eq:appendix-gauge-cokernel} is zero. If $b_1$ is another absolute
primitive, then $b_1-b_0$ is a global cocycle. Its restriction disappears in
the cokernel, so the remaining class does not depend on which primitive
was used.

The v0.29 exact sequence identifies this secondary cokernel with the kernel
of the map that forgets local gauge data from the relative obstruction
space. Exactness at that position does not make the forgetting map
surjective. It also does not make the localization map in
\eqref{eq:appendix-gauge-cokernel} surjective. These are separate arithmetic
questions. In particular, solving the absolute $H^2$ obstruction does not
by itself solve a family of prescribed local conditions.

\subsection{The reachable equation target is determined by syzygies}

Let $J:V\to W$ be a linear map over a field $k$. The correction equation
$Ju+e=0$ has a solution precisely when $e\in\operatorname{im}J$. A useful
description of this image uses all the dual relations among the equations.
Let
\[
 S=\ker(J^*:W^*\to V^*),\qquad
 C=\{w\in W:\phi(w)=0\text{ for every }\phi\in S\}.
\]
Every $Ju$ belongs to $C$. Conversely, if $w$ does not belong to
$\operatorname{im}J$, its nonzero class in $W/\operatorname{im}J$ is
detected by an algebraic linear functional. Pulling that functional back
to $W$ gives an element of $S$ that does not annihilate $w$. Hence
\begin{equation}
 \operatorname{im}J=C.
 \label{eq:appendix-compatible-target}
\end{equation}
This algebraic dual argument works without a finite-dimensional hypothesis.
It makes no assertion about continuous duals in a different topology.

Corestricting $J$ to $C$ makes it surjective onto its actual reachable
target. The equations have not changed. By contrast, adjoining variables
that absorb arbitrary errors changes the system unless their removal is
proved. A solution of an augmented equation is not automatically a solution
of the original equation. The original target and the quotient that records
unreachable errors must be kept separate.

\begin{example}[A reachable target with a visible failed correction]
Let $V=k^2$, $W=k^3$, and
$J(u,v)=(u,v,u+v)$. Its image is the plane $w_3=w_1+w_2$. The dual
kernel is generated by the functional
$\phi(w_1,w_2,w_3)=w_3-w_1-w_2$. The error $(1,2,3)$ is compatible
in every characteristic and is cancelled by $(u,v)=(-1,-2)$. The error
$(0,0,1)$ is incompatible in every field, since its value under $\phi$
is $1$. No choice of the two correction variables can cancel it.

If a third variable $t$ is added so that the derivative becomes
$J'(u,v,t)=(u,v,u+v+t)$, the enlarged map is surjective. It cancels the
second error with $(u,v,t)=(0,0,-1)$. But the nonzero value of $t$ records
exactly the failed original relation. This enlarged solution cannot be
declared a solution of the original equations by forgetting $t$. One must
show that the additional variable is removable, or that the target error
already lies in the original compatible plane.
\end{example}

For a nonlinear system $F(x)=0$, a relation among the residual Jacobian
rows is only a first-order identity. Consider
\[
 F(x)=(x,x^2),\qquad J_x(h)=(h,2xh).
\]
At the origin the derivative has horizontal image. The functional selecting
the second coordinate is a residual syzygy, but it does not annihilate
$F(x)$ for general $x$. The exact relation is the moving relation
\[
 (-x,1)F(x)=0.
\]
Furthermore,
\[
 F(x+h)-F(x)-J_x(h)=(0,h^2).
\]
Modulo $8$, taking $x=0$ and $h=2$ gives remainder $(0,4)$, outside the
horizontal target. Keeping only the syzygy at the origin would therefore
lose a genuine nonlinear error even within the same residual fibre.

The graded correction argument retains an exact moving relation
$B_xF(x)=0$ and a residual exactness statement
\begin{equation}
 \operatorname{im}\overline J=\ker\overline B.
 \label{eq:appendix-residual-exactness}
\end{equation}
If $F(x)=\pi^r e$, the moving identity gives $\pi^r B_xe=0$. A faithful
leading layer allows this equality to imply
$\overline B\,\overline e=0$. Equation
\eqref{eq:appendix-residual-exactness} then supplies a residual correction
$\overline u$ with $\overline J\,\overline u=-\overline e$. Lift
$\overline u$ to $u$ and put $x'=x+\pi^r u$. The linear part of $F(x')$
is divisible by $\pi^{r+1}$. For polynomial systems the remainder contains
at least two copies of $\pi^r$, so it is divisible by $\pi^{2r}$ and hence
by $\pi^{r+1}$ when $r\ge1$. This produces the next layer.

Over a suitable truncated principal coefficient ring, iteration terminates
at $\pi^N=0$ and produces a root of the original system. The recovered
v0.25 modules derive the necessary faithful-layer property under their
stated principal residue and nilpotence hypotheses. They still require
the moving identity and exactness. Equality of selected dimensions alone
does not prove either one. A correct dimension count can help prove
exactness after an inclusion is established, but cannot replace that
inclusion.

The normalization by $\pi^r$ deserves care. Writing
$F(x)=\pi^r e$ need not determine $e$ uniquely in a ring with torsion.
If $e'$ is another choice, then $\pi^r(e-e')=0$. The faithful leading
layer used by the theorem ensures that $e-e'$ has zero residue below the
terminal cutoff. Thus the residue correction problem is independent of
the chosen normalized error. The same hypothesis converts
$\pi^r B_xe=0$ into the residue relation needed for exactness. These are
two concrete reasons why the coefficient-ring condition is part of the
theorem. Formal cancellation of $\pi^r$ in a ring where it is a zero
divisor would be invalid.

The moving relation is evaluated at the current approximate root, rather
than at the initial point alone. The residual class of that root is kept
fixed during every correction, so its residual derivative and residual
relation operator remain the same. The exact relation changes over the
coefficient ring as the root changes. This is how the iteration retains
the full nonlinear equations while continuing to solve one fixed
residual linear problem. The distinction between a fixed residual
operator and a fixed integral operator is the mechanism that permits
redundant equations without an assumed integral right inverse.

There is also an integral warning. In $R=\mathbb Z/4\mathbb Z$, the
submodule $C=2R$ contains $2$, and $2$ is divisible by $2$ in $R$. Yet
$2$ is not divisible by $2$ inside $C$, since $2C=0$. Thus an ambient
divisibility statement need not be an internal divisibility statement in a
compatible target. An integral splitting cannot be inferred from the
residual vector-space picture. The graded argument avoids silently
assuming such a splitting by working with the normalized residue of each
actual error.

\subsection{Ordered powers have a uniform ultrametric remainder}

Let $E$ be a complete normed ring with $\norm{1}=1$, a submultiplicative
norm, and the ultrametric inequality
\[
 \norm{x+y}\le\max(\norm{x},\norm{y}).
\]
Assume $\norm A\le1$. For a nonnegative integer $d$, define
\begin{align}
 D_{A,d}(h)&=\sum_{i=0}^{d-1}A^ihA^{d-1-i},
 \label{eq:appendix-ordered-derivative}\\
 R_{A,d}(h)&=(A+h)^d-A^d-D_{A,d}(h).
 \label{eq:appendix-ordered-remainder}
\end{align}
For $d=0$ the sum is empty and both maps are zero. The order of the
factors in \eqref{eq:appendix-ordered-derivative} matters. Replacing this
operator by $dA^{d-1}h$ requires $A$ and $h$ to commute. The replacement
is invalid for general matrices.

The recursive identities are
\begin{align*}
 D_{A,d+1}(h)&=D_{A,d}(h)A+A^dh,\\
 R_{A,d+1}(h)&=R_{A,d}(h)(A+h)+D_{A,d}(h)h.
\end{align*}
They follow by multiplying $(A+h)^d$ on the right by $A+h$ and collecting
the terms with zero, one, and at least two occurrences of $h$. This
derivation does not rearrange any noncommuting factors.

For degree three, the linear term is
\[
 D_{A,3}(h)=A^2h+AhA+hA^2,
\]
and the nonlinear remainder is
\[
 R_{A,3}(h)=Ah^2+hAh+h^2A+h^3.
\]
The middle word $hAh$ cannot be absorbed into $Ah^2$ in a matrix ring.
The remainder estimate keeps all four words and bounds each by its
norm, so it respects the ordered expansion. This also explains why the
estimate is not an assertion of commutative binomial calculus in
disguise. The scalar bound remains uniform because the norm of a finite
ultrametric sum is controlled by its largest summand, rather than by the
number of summands.

\begin{lemma}[Uniform remainder estimate]
For $0\le r\le1$ and $\norm h,\norm k\le r$,
\[
 \norm{D_{A,d}(h)}\le\norm h,
 \qquad
 \norm{R_{A,d}(h)-R_{A,d}(k)}\le r\norm{h-k}.
\]
In particular $\norm{R_{A,d}(h)}\le r\norm h$.
\end{lemma}

\begin{proof}
The recursive formula for $D$ shows its additivity. Each term of its
ordered sum has norm at most $\norm h$, because every power of $A$ has
norm at most one. A finite ultrametric sum has norm at most the largest
term, proving the first estimate without a factor $d$.

Induct on $d$ for the second estimate. The case $d=0$ is immediate. Write
$D_d=D_{A,d}$ and $R_d=R_{A,d}$. Subtracting the remainder recursions gives
\begin{align*}
 R_{d+1}(h)-R_{d+1}(k)
 ={}&(R_d(h)-R_d(k))(A+h)\\
 &+R_d(k)(h-k)\\
 &+(D_d(h)-D_d(k))h+D_d(k)(h-k).
\end{align*}
The first term has norm at most $r\norm{h-k}$ by induction and
$\norm{A+h}\le1$. Taking $h=k$ and the other argument zero in the
inductive estimate gives $\norm{R_d(k)}\le r\norm k\le r$, so the
second term has the same bound. Additivity gives
$\norm{D_d(h)-D_d(k)}\le\norm{h-k}$, which bounds the third term.
Finally $\norm{D_d(k)}\le\norm k\le r$ bounds the fourth term.
Ultrametricity bounds their sum by the same maximum. Since $R_{A,d}(0)=0$,
taking $k=0$ gives the final claim.
\end{proof}

For $E=\operatorname{Mat}_n(\mathbb Q_p)$ use the entrywise supremum norm
and normalize $\abs p_p=p^{-1}$. Matrix multiplication satisfies
\[
 \abs{(XY)_{ij}}_p
 =\left|\sum_tX_{it}Y_{tj}\right|_p
 \le\max_t\abs{X_{it}}_p\abs{Y_{tj}}_p
 \le\norm X\norm Y.
\]
Taking the maximum over the finite set of entries proves
submultiplicativity. The same ultrametric calculation removes a matrix
dimension factor. Completeness is inherited from the finite product of
copies of $\mathbb Q_p$. These are the hypotheses required by the actual
matrix application, rather than properties borrowed from a complex matrix
norm.

\subsection{An exact inverse turns the local image into one finite test}

Assume that $D=D_{A,d}$ has an additive two-sided inverse $Q$, and that
$\norm{Qz}\le K\norm z$, with $K\ge0$. Choose $0\le r\le1$ with $Kr<1$. For a
target $U$, set $\Delta=U-A^d$. The proposed correction map is
\begin{equation}
 \Phi(h)=Q(\Delta-R_{A,d}(h)).
 \label{eq:appendix-contraction-map}
\end{equation}
Both inverse identities will be used. $DQ=1$ turns a fixed point into a
root. $QD=1$ proves the necessary condition from an existing root.

\begin{theorem}[Exact image of a small power ball]
Under these hypotheses, there exists a unique $h$ with $\norm h\le r$
and $(A+h)^d=U$ if and only if $\norm{Q\Delta}\le r$.
\label{thm:appendix-local-image}
\end{theorem}

\begin{proof}
Suppose $\norm{Q\Delta}\le r$. On the closed $r$-ball, the remainder
estimate and ultrametricity give
\[
 \norm{\Phi(h)}
 \le\max\bigl(\norm{Q\Delta},K\norm{R_{A,d}(h)}\bigr)
 \le\max(r,Kr\norm h)\le r.
\]
For $h,k$ in that ball,
\[
 \norm{\Phi(h)-\Phi(k)}\le Kr\norm{h-k}.
\]
The ball is complete and $Kr<1$. The contraction theorem gives a unique
fixed point $h$. Apply $D$ to its equation to obtain
$Dh=\Delta-R_{A,d}(h)$, which is exactly $(A+h)^d=U$.

Conversely, a root in the ball satisfies $\Delta=Dh+R_{A,d}(h)$. Applying
$Q$ gives
\[
 Q\Delta=h+QR_{A,d}(h).
\]
Its norm is at most $r$, since the second summand has norm at most
$Kr\norm h\le r$. Any such root is a fixed point of
\eqref{eq:appendix-contraction-map}, so the contraction also proves
uniqueness among all roots in this ball.
\end{proof}

The proof gives a controlled approximation process. Start at $h_0=0$ and
iterate $h_{j+1}=\Phi(h_j)$. Then $h_1=Q\Delta$, and
\[
 \norm{h_{j+1}-h_j}\le (Kr)^j\norm{h_1}.
\]
For $t>j$, ultrametricity bounds $\norm{h_t-h_j}$ by the largest of
the intervening increments, which is at most $(Kr)^j\norm{h_1}$. The
limit therefore satisfies the same bound for its distance from $h_j$.
With $r=p^{-m}$ and $K=p^s$, every iteration gains at least $m-s$
digits in this estimate. This is a quantitative consequence after the
inverse and target test have been certified. It is not a search over
unrelated roots at successively deeper levels.

Uniqueness is restricted to the chosen ball. Two different balls can each
contain a root of the same target. A failed target congruence rules out
the specified residue prefix, while an accepted congruence determines
one root with that prefix. Covering all possible prefixes requires a
separate argument, especially near centers where the derivative is
singular or where the available precision is too shallow.

For integral matrices suppose the inverse has the certified form
\[
 Q=p^{-s}B,\qquad s\in\mathbb N,
\]
where $B$ is an integral linear operator on the matrix entries. In a fixed
ordering of the $n^2$ entries it is an $n^2$ by $n^2$ matrix over
$\mathbb Z_p$. This is an operator on the matrix space, not ordinary
left multiplication by an $n$ by $n$ matrix. Integral coefficients imply
$\norm{Bz}\le\norm z$, so $K=p^s$ is a valid bound.

There is a finite exact certificate for the inverse. For every matrix unit
$E_{ij}$ check
\[
 QD(E_{ij})=E_{ij},\qquad DQ(E_{ij})=E_{ij}.
\]
Linearity and the matrix-unit basis extend both identities to every matrix
over $\mathbb Q_p$. These equalities are exact in the coefficient field.
They are not merely congruences modulo $p$. In particular, an operator
singular over $\mathbb Q_p$ cannot satisfy them.

Take $r=p^{-m}$ with $m>s$. Then $Kr=p^{s-m}<1$, and
\begin{align}
 \norm{Q(U-A^d)}\le p^{-m}
 &\Longleftrightarrow \norm{B(U-A^d)}\le p^{-(m+s)}\notag\\
 &\Longleftrightarrow
 B\operatorname{vec}(U-A^d)=0\pmod {p^{m+s}}.
 \label{eq:appendix-finite-criterion}
\end{align}
The last congruence is entrywise. It depends only on the reductions of the
target, center, and integral operator modulo $p^{m+s}$. The theorem
therefore decides the existence and uniqueness of a root in
$A+p^m\operatorname{Mat}_n(\mathbb Z_p)$ by one finite target test, once
the exact inverse certificate is supplied.

The root constructed over $\mathbb Q_p$ is integral. Its center is
integral, and every correction entry has norm at most $p^{-m}\le1$.
If $d>0$ and $U$ is a unit, the root $F$ is also a unit. Indeed, $F$
commutes with $U=F^d$ and hence with $U^{-1}$, and
$F^{d-1}U^{-1}$ is a two-sided inverse of $F$. Thus the unit version has
an actual output in $\operatorname{GL}_n(\mathbb Z_p)$.

The value of $s$ measures denominators of the whole inverse derivative.
It is not automatically $v_p(d)$. The proof makes no statement about
shallow balls with $m\le s$, where its contraction estimate is insufficient.
It also does not provide automatic certificate discovery for arbitrary
$p$-adic input. Exact rational examples can supply certificates by ordinary
linear algebra, but a general certified implementation is an additional
task.

\subsection{A noncommuting unipotent square example}

Work over $\mathbb Z_2$ and take
\[
 A=\begin{pmatrix}1&1\\0&1\end{pmatrix},\qquad d=2.
\]
For $X=\begin{pmatrix}a&b\\c&e\end{pmatrix}$, the ordered derivative is
\[
 D_{A,2}(X)=AX+XA
 =\begin{pmatrix}2a+c&a+2b+e\\2c&c+2e\end{pmatrix}.
\]
Using the entry order $(a,b,c,e)$, its matrix and a scaled inverse are
\[
 J=\begin{pmatrix}
 2&0&1&0\\1&2&0&1\\0&0&2&0\\0&0&1&2
 \end{pmatrix},\qquad
 B=\begin{pmatrix}
 2&0&-1&0\\-1&2&1&-1\\0&0&2&0\\0&0&-1&2
 \end{pmatrix}.
\]
Direct multiplication gives $BJ=JB=4I_4$. Therefore $Q=\tfrac14 B$
is the exact inverse and $s=2$. The entry $\tfrac14$ in this inverse
prevents clearing its denominators with only one factor of $2$.

The residual derivative is singular. In characteristic two,
$D_{A,2}(I_2)=2A=0$, while $I_2\ne0$. Nevertheless $J$ is invertible
over $\mathbb Q_2$, with determinant $16$. Residual singularity and
singularity over the fraction field are different tests.

Choose $m=3$, so the local root ball is $A+8\operatorname{Mat}_2(\mathbb Z_2)$
and the finite target test is modulo $2^{m+s}=32$. Consider
\[
 F=\begin{pmatrix}1&1\\8&1\end{pmatrix},\qquad
 U=F^2=\begin{pmatrix}9&2\\16&9\end{pmatrix}.
\]
Here
\[
 AF=\begin{pmatrix}9&2\\8&1\end{pmatrix},\qquad
 FA=\begin{pmatrix}1&2\\8&9\end{pmatrix},
\]
so the root does not commute with the center. Since
$A^2=\begin{pmatrix}1&2\\0&1\end{pmatrix}$, the target difference has
entry vector $(8,0,16,8)$. Applying $B$ gives
\[
 B\operatorname{vec}(U-A^2)=(0,0,32,0).
\]
The congruence modulo $32$ passes. The theorem supplies uniqueness in the
specified ball, beyond the direct multiplication that exhibits this root.

For a negative control, take
\[
 U_{\mathrm{bad}}=\begin{pmatrix}1&10\\0&1\end{pmatrix}.
\]
The difference from $A^2$ is $(0,8,0,0)$ in entry coordinates, and $B$
sends it to $(0,16,0,0)$. This is nonzero modulo $32$, so no square root
of $U_{\mathrm{bad}}$ lies in that ball. This is a failure for the
specified prefix. It is not a claim that the target has no square root
anywhere in a larger coefficient field.

The recovered exact checker compares the finite image with the finite
criterion. Its unipotent square case tests $256$ corrections at the
stated horizon and $4096$ corrections one level deeper. Both give the same
$16$ target residues as the criterion. These are bounded corroborating
calculations. Infinite existence and uniqueness follow from the general
contraction proof and its exact hypotheses.

There is a genuinely singular contrast. Let
$A_0=\operatorname{diag}(1,-1)$ over $\mathbb Q_p$. Then
$D_{A_0,2}(E_{12})=0$, although $E_{12}\ne0$. Thus this derivative has
no inverse over the fraction field, including when $p$ is odd and $2$ is
a unit. The certificate theorem cannot cover this center. A universal
claim that invertible centers always have invertible power derivatives
would be false.

\subsection{Finite levels give a coherent branch only for a joint problem}

Let $X_r$ be finite sets with genuine reduction maps
$\pi_{r,t}:X_t\to X_r$ for $r\le t$. Assume identity and composition
laws for these maps. Let $P_r\subseteq X_r$ be solution sets, and assume
reduction preserves solutions. Then
\begin{equation}
 \left(\forall r,\ P_r\ne\varnothing\right)
 \Longleftrightarrow
 \left(\exists (x_r)_r,\ x_r\in P_r\text{ and }
 \pi_{r,t}(x_t)=x_r\text{ for }r\le t\right).
 \label{eq:appendix-branch-criterion}
\end{equation}
The nontrivial implication is finite branch compactness. For each depth,
a solution at that depth gives a finite path of solutions at all earlier
depths. There are arbitrarily long paths. The resulting tree has finite
branching because every level is finite. An infinite branch follows from
the finite-branch form of K\"onig's lemma. The reduction maps on the
solution sets need not be surjective. Compactness selects persistent
choices instead of extending a previously chosen node at every step.

This matters for simultaneous matrix prescriptions. At level $r$, a node
can contain a shared matrix $F_r$ together with one local conjugator
$z_{i,r}$ for each member of a finite family. The equations are
\[
 z_{i,r}F_r^{d_i}z_{i,r}^{-1}=U_i\pmod {p^r},
\]
with the prescribed prefix of $F_r$ and prescribed residual identity for
each conjugator. These are a single joint predicate on a finite tuple.
The arithmetic equations and prefixes are preserved by reduction.
Compactness therefore produces compatible towers for the shared matrix
and for all local conjugators, provided this joint problem is nonempty at
every depth.

An individual solution for each prescription is weaker. Even the scalar
requirements $F=1$ and $F=-1$ modulo an odd prime have solutions
separately, but have no common solution. Likewise, solving every member
of a family at each depth with a different shared matrix does not establish
the premise of \eqref{eq:appendix-branch-criterion}. The compatible branch
must come from one family of jointly satisfiable nodes.

Nor does a successful shallow node necessarily extend. For example, the
scalar prefix $F=4$ modulo $9$ satisfies $F^3=1$ there. Its lifts modulo
$27$ are $4,13,22$, all with cube $10$. That branch stops at the next
level. Merely checking a finite number of successful levels gives no
all-depth conclusion without another theorem.

Conversely, for a downward-closed tower of finite solution sets indexed
by all natural-number levels, absence of an infinite branch implies that
some finite level is empty. This gives a finite
obstruction witness in principle. It supplies no uniform computable bound
on the depth of that witness. The matrix inverse theorem above provides
an explicit horizon only in its certified nonsingular sector with $m>s$.
It removes the all-depth premise for one local power problem in that sector.
It does not remove the joint all-depth premise for arbitrary simultaneous
prescriptions.

Finally, compatible integral matrix residues determine an actual matrix
over $\mathbb Z_p$ entry by entry. Compatible inverse residues determine
its inverse when the nodes are units. Equality of all reductions then
implies equality of the integral matrices. These reconstruction steps turn
the selected branch into a genuine integral solution. The remaining
arithmetic attachment is not part of that reconstruction.

\subsection{What the formal receipts establish}

The recovered archives preserve the source statements and historical
verification receipts. The v0.38 receipt records a successful build of its
seven new modules against the pinned Mathlib environment. The earlier
interfaces described here are drawn from their recovered source files and
their recorded histories. Archive authentication verifies the recovered
bytes. It is not a new compilation of the mathematical statements.

A fresh Lean replay was attempted during preparation of this edition, but
the portable kit could not be recovered because its download returned
HTTP 403. No fresh compiler invocation was therefore completed for this
edition. The accompanying coverage report records that failure and
distinguishes it from the historical successful receipt. The new exact
finite controls also have their own tested domains. They do not substitute
for a kernel replay or extend a conditional theorem to a missing
arithmetic hypothesis.

The formal program supplies explicit cocycle quotients, exact obstruction
criteria, careful nonlinear correction interfaces, compatible tower
reconstruction, and an effective finite criterion for a substantial local
matrix sector. To reach Strong Artin, these mechanisms still need the
arithmetic identifications and global existence arguments stated in the
research boundary. Their mathematical value comes from making those
remaining premises visible, while proving the local implications in full.
